Recent advances in large language models have enabled increasingly capable AI assistants, yet deploying these agents in real-time communication platforms like Discord presents unique challenges. Users expect assistants to be helpful and proactive, but also safe and accurate. Unproactive agents miss opportunities to assist, while overly aggressive agents can annoy users or perform unwanted actions. Furthermore, language models are prone to hallucinations and may lack up-to-date or domain-specific knowledge needed for effective assistance in specialized communities.

Existing approaches to AI agents in communication platforms typically fall into two categories: reactive bots that only respond to explicit commands, or proactive agents that initiate actions without sufficient user oversight. Reactive agents fail to anticipate user needs, while proactive agents without oversight can compromise safety and accuracy. Recent work on tool use in language models (Schick et al., 2023; Yao et al., 2023) shows promise for extending model capabilities, but often lacks the safety mechanisms needed for real-world deployment. Similarly, retrieval-augmented generation approaches improve factual accuracy but may not address the proactive assistance and safety concerns inherent in interactive environments.

Our work builds upon and extends several lines of research. First, we contribute to the growing literature on proactive AI assistants (Horvitz, 1999; Early et al., 2005; Baldauf et al., 2007) by applying these principles to the Discord platform and combining them with safety mechanisms. Second, our approval workflow builds on research in human-in-the-loop machine learning (Amershi et al., 2014; Kulesza et al., 2015) and shared control systems (Dragan & Srinivasa, 2013; Hayes & Scassellati, 2015). Third, our integration with external reasoning engines connects to recent advances in retrieval-augmented generation for large language models (Lewis et al., 2020; Guu et al., 2020; Izacard & Grave, 2021), particularly noting how systems like NotebookLM implement these techniques for grounded reasoning. Finally, our work on reducing hallucinations through knowledge grounding relates to recent efforts to detect and mitigate factual errors in language model outputs (Ji et al., 2022; Zhang et al., 2023; Huang et al., 2023).

To our knowledge, we are the first to combine proactive behavior, structured approval workflows, and external reasoning integration in a single framework for communication platform agents. While individual components have been explored previously, the specific combination and evaluation in the context of Discord-based AI assistants represents a novel contribution.